\section{Variational state preparation}
\label{sec:ansatz}

\subsection{Optimization targets}
\label{sec:points}

Variational optimization targets each point of the two scan lines, $\delta=-2$ and $\delta=0.85$, on the $99$-point $s$ grid ($s=0.01$--$0.99$): at every grid point the three initial states are optimized independently and the lowest of the three minima is reported, so the target Hamiltonians sample the trivial, topological, and AFM regions mapped in Fig.~\ref{fig:D04}(d). The lowest-energy initial-state branches along the two lines are shown by the background bands in Fig.~\ref{fig:D08D09}(a,b).

\subsection{Variational circuit structure}

Our variational circuits follow the variational eigenvalue solver
paradigm~\cite{peruzzo2014variational}: prepare an initial state, then
apply a parameterized ansatz. The classical loop minimizes
\begin{equation}
E(\theta)=\langle\psi(\theta)|H|\psi(\theta)\rangle,
\end{equation}
where $|\psi(\theta)\rangle$ is the circuit output at parameters
$\theta$. At each target point, the three reference states are optimized independently at every depth $p=1,2,3$.
We screen random parameter vectors together with the zero vector, refine the low-energy candidates with local L-BFGS, and retain the lowest-energy converged result in each branch.
The minimum across the three branches determines $E^*$ and its parameters.
The symmetry-preserving circuit retains the initial symmetry quantum numbers.
The scan runs point by point along the $\delta=-2$ and $\delta=0.85$ lines. Variational algorithms of this form are surveyed in
Refs.~\cite{cerezo2021variational,tilly2022variational}, which collect
cost-function design, ansatz families, and best practices for VQE.

The ideal-state optimization is performed classically; the processor executes the resulting circuits to measure correlators.
For a normalized trial state in a sector with a unique ground state, spectral decomposition gives
\begin{equation}
1-|\langle\psi_0|\psi(\theta)\rangle|^2
\leq \frac{E(\theta)-E_0}{\Delta_{\rm sec}}.
\label{eq:energy-fidelity}
\end{equation}
Equation~\eqref{eq:energy-fidelity} relates the trial-state overlap to the energy residual through the sector gap $\Delta_{\rm sec}$.

Our bond rotations form a Hamiltonian variational ansatz: the generators follow the interaction terms and their angles are optimized independently~\cite{wecker2015progress}.
The XXZ setting has been studied within this family~\cite{wiersema2020exploring}.
Building symmetry into state-preparation circuits restricts the search to relevant sectors~\cite{gard2020efficient}.
Different reference states supply different initial entanglement patterns, a relevant resource in variational preparation of nontrivial states~\cite{ho2019efficient}.
The three reference states share the same finite-chain quantum numbers and provide distinct entanglement structures; their optimized energies determine the branch selected at each parameter point.

Trainability also depends on circuit and cost structure.
Random circuits approaching unitary designs can exhibit exponentially small gradients~\cite{mcclean2018barren}, and local versus global costs can behave differently even at shallow depth~\cite{cerezo2021cost}.
Identity-block initialization~\cite{grant2019initialization} and model-based circuit design~\cite{wiersema2020exploring} motivate structured starting points.
Depth-dependent optimization transitions~\cite{campos2021abrupt} and noise-induced gradient suppression under specified noise models~\cite{wang2021noise} further characterize the roles of circuit depth and noise in variational optimization.

Let $M=L/2$. The three reference initial states lie in the $P=-2$ sector and are shown in Fig.~\ref{fig:D06}. They are paired with an
orbit ansatz in which mirror-symmetric bond pairs share rotation angles
(odd bonds $O_j\leftrightarrow O_{M+1-j}$, even bonds
$E_j\leftrightarrow E_{M-j}$). The ansatz is built in two steps:
first fix the entanglement structure with a phase-matched initial state,
then dress all bonds with rotations tied by the reflection symmetry, so
the parameter count grows with the number of bond orbits rather than
bonds. Each sublayer evolves under ($\alpha$ labels the orbit,
$\langle kl\rangle$ a bond):
\begin{equation}
U_{o}=\prod_{\alpha}\prod_{\langle kl\rangle\in\mathcal{O}_{\alpha}}e^{-i\theta_{o,\alpha}(X_{k}X_{l}+Y_{k}Y_{l})/2}e^{-i\varphi_{o,\alpha}Z_{k}Z_{l}/2},
\end{equation}
\begin{equation}
U_{e}=\prod_{\alpha}\prod_{\langle kl\rangle\in\mathcal{E}_{\alpha}}e^{-i\theta_{e,\alpha}(X_{k}X_{l}+Y_{k}Y_{l})/2}e^{-i\varphi_{e,\alpha}Z_{k}Z_{l}/2}.
\end{equation}
On the same bond $XX{+}YY$ commutes with $ZZ$, and disjoint bonds within
a sublayer commute; the odd and even sublayers do not commute with each
other, so both are kept. Each bond is implemented as an
$R_{XX}(\theta)R_{YY}(\theta)R_{ZZ}(\varphi)$ gate sequence. Each orbit
carries one $(\theta,\varphi)$ pair, with $L$ parameters per layer for the $L=8$ chain considered here (halved
at $\delta=0$ where $\theta=\varphi$); the ansatz preserves
$Z_{\rm tot}$, $\prod X$, and $R$, staying in the $P=-2$ sector. The
sublayer overlapping the initial-state pairing acts last: even-first
for trivial and AFM initial states, odd-first for topological. One schematic per phase is shown in Fig.~\ref{fig:D06}(a)--(c): initial state
plus required ansatz circuit for the trivial, topological, and
AFM cases.
For ideal circuits, the trial states remain in the target symmetry sector. Section~\ref{sec:comparison} compares measured and simulated correlators.
Figure~\ref{fig:D08D09}(a,b) shows the $p=1$ cost functions
against the exact ground-state energy.

\begin{figure*}[htbp]
\centering
\includegraphics[width=\textwidth]{{../../figs/exp04/exp04_D06}}
\caption{Initial states and required ansatz circuits for the three phases:
(a) trivial; (b) topological; (c) AFM. The dashed box $U^{(k)}$ is the $k$th variational layer;
$\theta^{(k)}_j$ ($\varphi^{(k)}_j$) is the $XX{+}YY$ ($ZZ$) angle shared by
the $j$th mirror-symmetric bond orbit in layer $k$ (odd
$O_j\leftrightarrow O_{M+1-j}$, even $E_j\leftrightarrow E_{M-j}$;
$\theta=\varphi$ at $\delta=0$). Even-bond sublayer first in (a) and (c),
odd-bond sublayer first in (b). \label{fig:D06}}
\end{figure*}

\begin{figure*}[htbp]
\centering
\includegraphics[width=\textwidth]{{../../figs/exp06/exp06_D08D09}}
\caption{Cost functions at $p=1$ and the experiment--simulation comparison.
(a),(b) Cost functions at $p=1$ versus $s$ for (a) $\delta=-2$ and (b) $\delta=0.85$:
$E_0$ (exact, blue solid) with the three initial-state costs
(triv green dashed, topo orange dotted, AFM black dashdot);
background bands identify the initial-state branch minimizing the $p=1$ cost at each grid point; circles identify the branch with the lowest mean optimized cost within each interval.
(c),(d) Combined quantity $Q$ for (c) $\delta=-2$ and (d) $\delta=0.85$:
simulation $p=1$ (black) and $p=2$ (gray) fixed-scale curves with
experimental window scatter ($p=1$, 16-unit window of the 500-unit round; colors indicate
the lowest-energy initial-state branches in (a),(b)).
\label{fig:D08D09}}
\end{figure*}
